\documentclass[11pt]{article}
\usepackage[margin=1in]{geometry}
\usepackage[T1]{fontenc}
\usepackage{lmodern}
\usepackage{amsmath,amssymb}
\usepackage{booktabs}
\usepackage[round,authoryear]{natbib}
\usepackage{hyperref}
\usepackage{xurl}
\usepackage{microtype}
\hypersetup{
  hidelinks,
  pdftitle={A Replication and Fragility Note on Credit Growth and Downturn Severity},
  pdfauthor={Akufor Agho Fru}
}
\title{A Replication and Fragility Note on\\Credit Growth and Downturn Severity}
\author{Akufor Agho Fru}
\date{}

\begin{document}
\maketitle

\begin{abstract}
\citet{jst2013} report that recessions that follow credit-intensive expansions are deeper and
have slower recoveries. I re-estimate their aggregate Table 8 in Python. All 20 coefficients
and standard errors match the authors' archived output at three decimals, with 121 events.
On release 6 of the Jord\`a-Schularick-Taylor database, the relation between pre-downturn
credit growth and downturn severity is imprecise and moves with data vintage, episode
dating, and controls. On the 116 events common to both vintages, the excess-credit
coefficients move toward zero at horizons one to three, but none changes sign. With release
6 outcomes and controls, the five-year cumulative-loss coefficient is $-1.097$ (SE 0.446)
under an annual-decline rule, $-0.013$ (SE 0.432) under the JST chronology, and $-0.401$
(SE 0.501) on dates common to both rules. A separate design regresses nine severity
outcomes on three-year credit growth with house-price, other macro, country, and decade
controls, not the JST design. On 144 to 178 complete-case episodes, every Holm-adjusted
$p$-value is 1.000. By construction, a 0.1 percent annual output cost
and the observed frequency of 0.114 downturns per country-year imply that restraint must
avoid 0.87 GDP-percent-years of loss per downturn to break even. I do not express this as a
share of measured loss. The loss measure floors yearly gaps at zero, counts overlapping
windows, and includes war-year downturns, and it implies an annual trend loss of 4.5 to 11.9
percent of GDP across corrections. The note does not estimate how any policy changes
downturn losses.
\end{abstract}

\section{Introduction}
Credit booms are a central concern of macroprudential policy because they predict financial
crises. \citet{schularicktaylor2012} show that lagged credit growth predicts financial crises
in 14 advanced economies from 1870 to 2008. \citet{gourinchasobstfeld2012} find that domestic
credit expansion is among the most robust crisis predictors in advanced and emerging
economies after 1973. \citet{greenwood2022} show that rapid three-year growth in credit and
asset prices raises the probability of a crisis within three years to about 40 percent.
\citet{msv2017} show that a rise in household debt to GDP predicts lower output growth over
the medium run in 30 countries. \citet{sufitaylor2022} survey this evidence.

A second claim concerns severity, not frequency. \citet{jst2013} use local projections on 14
countries from 1870 to 2008. They find that recessions that follow larger credit build-ups
are deeper and have slower recoveries, in normal recessions as well as in financial ones.
\citet{claessens2012} find that recessions with financial disruptions, in particular house
price busts, are longer and deeper. \citet{bridges2017} study 130 downturns in 26 advanced
economies since the 1970s. They find that rapid credit growth before a recession predicts a
deeper and longer downturn, with or without a banking crisis, and that credit growth
predicts severity better than the level of debt. This evidence supports the severity link.
The imprecise release 6 estimates below do not contradict it, because they use a different
sample, period, and design.

The literature already shows that this severity link depends on the type of credit, the
features of the boom, and the dating of events. On credit type, \citet{jst2016mortgaging}
find that real estate lending booms are followed by deeper recessions and slower
recoveries. \citet{jkst2022} find that corporate debt booms do not show the same pattern,
except when debt resolution is slow. \citet{mullerverner2024} find that credit to the
non-tradable sector predicts slowdowns, while credit to the tradable sector does not.
\citet{jst2016sovereigns} find that the risk comes mostly from private, not public, debt.
On boom features, \citet{jst2015bubbles} find that asset price bubbles are most damaging
when credit finances them, and \citet{krishnamurthymuir2025} find that crisis severity
depends on credit losses interacted with prior credit growth. \citet{rsw2021} find that
credit booms with house price booms and a rising loan-to-deposit ratio are much more likely
to end in a systemic banking crisis than other credit booms. On dating,
\citet{romerromer2017} find that the output cost of crises in advanced countries is moderate
and sensitive to specification and sample. \citet{baron2021} show that bank equity declines
identify banking crises that narrative chronologies miss. \citet{bordohaubrich2017} find
that US recessions with financial crises are generally followed by fast recoveries, with
three exceptions: the 1930s, the early 1990s, and the recovery after 2008. Their answer
depends on how crises are defined. This note does
not add a new dimension of fragility. It checks the aggregate Table 8 estimate itself.

Cost-benefit work on boom restraint already exists. \citet{svensson2017} compares the costs
and benefits of leaning against the wind with monetary policy. He uses estimated effects of
the policy rate on the probability and size of crises, and he finds that costs exceed
benefits for representative estimates. \citet{bgnn2020} estimate how macroprudential and
monetary tightening shift the distribution of future output growth. They then compute net
benefits with a loss function. Both papers do cost-benefit analysis with estimated policy
effects. This note does not do that. It does not estimate the effect $D(m)$ of a policy $m$
on downturn loss.

This note asks a narrower question. Is the credit-severity association in \citet{jst2013}
stable when the data are revised and when the episode definition changes? It also asks how
large an avoided loss a costly restraint would need to break even. The contributions are
these:
\begin{enumerate}
\item A check of the aggregate \citet{jst2013} Table 8 estimate. An independent Python
  implementation matches the authors' archived output at three decimals (RQ1).
\item A fragility analysis on release 6 data. It holds the original design fixed and varies
  the data vintage, the episode dating rule, and the control set one at a time (RQ2).
\item A conditional break-even calculation. It converts an annual output cost into the
  avoided loss per downturn that the cost requires. It also shows why that loss should not
  be expressed as a share of the measured loss per downturn (RQ3).
\end{enumerate}
\textbf{RQ1.} Does an independent implementation reproduce Table 8 of \citet{jst2013}?
\textbf{RQ2.} On release 6 data, how sensitive is the credit-severity association to
vintage, dating, and controls? \textbf{RQ3.} Given the observed downturn frequency and loss
distribution, what avoided loss per downturn would a stated annual cost require?

The note does not overturn \citet{jst2013}. Their archived estimates are matched. The finding
is that the aggregate association is imprecise on revised data and depends on analyst
choices, which is consistent with the literature above. The note makes no causal or policy
claim. I found no published replication of \citet{jst2013} in the Institute for Replication
reports and discussion papers or in the FORRT replication atlas. ReplicationWiki was
unavailable when I checked. The nearest published replication I found is
\citet{summers2017}, which replicates \citet{schularicktaylor2012}. I make no priority claim.

\section{Data}
The reproduction uses the authors' replication archive \citep{jstcode}. It contains the
original panel (\texttt{private.dta}), the Stata programs, and the table output
\texttt{table8.tex}. The fragility analysis uses release 6 of the Jord\`a-Schularick-Taylor
Macrohistory Database \citep{jstdatabase}, available at
\url{https://www.macrohistory.net/database/}. I downloaded the file on 4 November 2026 and
confirmed on 18 November 2026 that the same URL served a file with the same SHA-256 hash. Release 6 has 2,718 country-year observations for
18 advanced economies from 1870 to 2020. The main activity series is real GDP per capita.
Credit is total loans to the non-financial private sector (\texttt{tloans}).

Release 6 \texttt{tloans} is the closest available series to the original \texttt{loans1}.
I could not confirm that the two definitions are identical. A vintage comparison therefore
mixes data revision with any definition change.

\section{RQ1: Reproduction of Table 8}
Table 8 of \citet{jst2013} regresses the cumulative change in log real GDP per capita at
horizons one to five on normal and financial recession indicators and on their interactions
with excess credit. Excess credit is the change in the credit-to-GDP ratio per year of the
preceding expansion. The model includes seven year-zero controls, their first lags, and
country effects. I translated the Stata program to Python with classical OLS standard errors.

All 20 coefficients and standard errors match the archive's \texttt{table8.tex} at three
decimals. Each horizon uses 121 events, 92 normal and 29 financial. The normal-recession
excess-credit coefficients are $-0.264$, $-0.684$, $-0.771$, $-0.932$, and $-0.711$ at
horizons one to five. The financial-recession coefficients are $-0.401$, $-0.991$,
$-0.374$, $-1.298$, and $-0.895$. I compared against the table file in the authors' archive,
not against a scan of the journal page. I did not run Stata.

\section{RQ2: Fragility on release 6 data}
\subsection{Data vintage on a fixed event set}
I hold the original event dates, recession labels, expansion durations, controls, lags,
fixed effects, and horizons fixed, and replace the series with release 6. Five of the 121
events lack complete release 6 inputs. Table~\ref{tab:vintage} uses the 116 common events.
At horizons one to three, all six release 6 credit interactions are closer to zero than the
original ones. At horizon two, the normal coefficient falls from $-0.645$ to $-0.290$ and the
financial coefficient falls from $-0.980$ to $-0.315$. At horizons four and five, the
financial coefficient stays large in magnitude, but its standard error rises. Release 6
standard errors for the financial interaction are larger at every horizon. No coefficient
changes sign. On all 148 release 6 events with complete inputs on the original dates, all
ten excess-credit coefficients are also negative, from $-0.177$ to $-1.022$. The sign
changes reported below come only from the smaller subsample, the control set, or the
episode rule. I did not run a formal test of the coefficient difference between vintages.
The two estimates use different data on the same events, so the standard errors in the
table do not support a simple difference test.

\begin{table}[htbp]
\centering
\small
\caption{Excess-credit coefficients of the Table 8 design on the 116 events common to the
original data and release 6. Classical OLS standard errors are in parentheses. Each model
has 88 normal and 28 financial recessions.}
\label{tab:vintage}
\begin{tabular}{lrrrr}
\toprule
Horizon & Normal, original & Normal, R6 & Financial, original & Financial, R6 \\
\midrule
1 & $-0.300$ (0.171) & $-0.188$ (0.178) & $-0.405$ (0.217) & $-0.247$ (0.253) \\
2 & $-0.645$ (0.299) & $-0.290$ (0.314) & $-0.980$ (0.378) & $-0.315$ (0.447) \\
3 & $-0.788$ (0.411) & $-0.276$ (0.408) & $-0.420$ (0.520) & $-0.250$ (0.580) \\
4 & $-0.860$ (0.526) & $-0.402$ (0.543) & $-1.296$ (0.664) & $-1.042$ (0.772) \\
5 & $-0.648$ (0.563) & $-0.415$ (0.554) & $-0.901$ (0.711) & $-1.155$ (0.789) \\
\bottomrule
\end{tabular}
\end{table}

Using all 148 events changes the sample as well as the data, so Table~\ref{tab:vintage}
uses only the common 116.

\subsection{Episode dating}
The second test changes the episode rule. The annual-decline rule starts an episode at a
local output peak when real GDP per capita falls in the next year. The JST chronology uses
the Bry-Boschan peaks from the original archive. Both samples use the same release 6
outcomes, seven pre-peak predictors, country and decade effects, and country-clustered
standard errors. The predictors are three-year credit growth, credit to GDP, house-price
growth, investment share, prior real growth, inflation, and the short rate. In this note,
three-year credit growth is the change in private credit to GDP over the three years before
the peak, in percentage points. Inference uses a
$t$ distribution with one fewer degree of freedom than the number of clusters.

In the shared coverage of 14 countries from 1870 to 2002, the annual rule finds 225 events
and the JST chronology has 223 peaks. Of these, 182 have the same country and year.
Table~\ref{tab:rules} reports the credit-growth coefficient for two outcomes. The five-year
cumulative-loss coefficient is $-1.097$ (SE 0.446) under the annual rule, $-0.013$ (SE 0.432)
under the JST chronology, and $-0.401$ (SE 0.501) on the matched dates. Only the annual-rule
estimate has a 95 percent interval that excludes zero ($-2.06$ to $-0.13$, raw $p=0.029$).
After a Holm adjustment across the eight tests in the table, its $p$-value is 0.229. The
two full-coverage annual-rule tests also appear in the nine-outcome family of
Table~\ref{tab:severity}, so they enter two Holm families. The
sign is also the opposite of the original finding, because a negative coefficient means
smaller losses after faster credit growth. Which events enter the sample changes the
estimate. This comparison does not separate the dating rule from sample composition.

\begin{table}[htbp]
\centering
\small
\caption{Coefficient on three-year credit growth under two episode rules, release 6 data.
Country-clustered standard errors are in parentheses. All models include the six other
predictors and country and decade effects. Outcomes are in percent and GDP-percent-years.}
\label{tab:rules}
\begin{tabular}{lrrr}
\toprule
Event sample & Peak-to-trough & Five-year loss & Events \\
\midrule
Annual rule, full R6 coverage & 0.093 (0.105) & $-0.005$ (0.441) & 178 \\
Annual rule, shared coverage & $-0.076$ (0.158) & $-1.097$ (0.446) & 137 \\
JST chronology, shared coverage & 0.031 (0.095) & $-0.013$ (0.432) & 126 \\
Matched dates & $-0.011$ (0.124) & $-0.401$ (0.501) & 109 \\
\bottomrule
\end{tabular}
\end{table}

\subsection{Controls on a fixed event set}
The third test holds the event set fixed and changes the specification one step at a time.
Of the 116 common events, 21 lack house-price growth and one lacks three-year credit growth.
That leaves 94 events in 14 countries. The base is the release 6 Table 8 design with
five-year GDP growth as the outcome. On these 94 events, the five-year excess-credit
coefficients are 0.062 for normal recessions and 0.037 for financial recessions. On the 116
events they were $-0.415$ and $-1.155$. The sample change alone moves both coefficients from
negative to near zero.

\begin{table}[htbp]
\centering
\small
\caption{Specification bridge on the fixed 94-event release 6 sample. Coefficients are per
unit of the predictor. Country-clustered standard errors are in parentheses. Each row
changes one element relative to the row above. In the first six rows the outcome is
five-year GDP growth, so a negative coefficient means weaker growth after more credit. In
the last two rows the outcome is a loss, so the sign convention flips: a positive
coefficient means a larger loss.}
\label{tab:bridge}
\begin{tabular}{p{4.2cm}p{2.6cm}rrr}
\toprule
Model & Predictor & Normal & Financial & Common \\
\midrule
R6 Table 8 design & Excess credit & 0.062 (0.391) & 0.037 (1.159) & \\
Switch predictor & 3-year credit & 0.118 (0.170) & 0.169 (0.312) & \\
Current controls & 3-year credit & $-0.217$ (0.092) & $-0.135$ (0.132) & \\
Add decade effects & 3-year credit & 0.030 (0.153) & $-0.127$ (0.335) & \\
Common slope & 3-year credit & & & $-0.001$ (0.184) \\
Drop event-type intercepts & 3-year credit & & & $-0.122$ (0.188) \\
Outcome: peak-to-trough & 3-year credit & & & 0.028 (0.124) \\
Outcome: five-year loss & 3-year credit & & & 0.230 (0.427) \\
\bottomrule
\end{tabular}
\end{table}

Table~\ref{tab:bridge} shows one interval that excludes zero. Under the current controls
without decade effects, the normal-recession coefficient is $-0.217$ (95 percent interval
$-0.416$ to $-0.018$). Adding decade effects moves it to 0.030. The bridge reports 12
coefficients without a multiple-testing adjustment, so a single interval that excludes zero
is weak evidence.

\subsection{All downturns on release 6}
The last test applies the annual-decline rule to the full release 6 panel. It finds 299
downturns in 18 countries. Of these, 59 overlap a systemic banking crisis year. The median
peak-to-trough decline is 3.11 percent, the mean is 6.27 percent, and the 90th percentile is
14.79 percent. The median recovery to the prior peak takes three years.

Each of nine severity outcomes is regressed on the seven pre-peak predictors with country
and decade effects. The outcomes are peak-to-trough decline, five- and ten-year cumulative
loss against a pre-episode linear trend, the same against a flat prior-peak path, recovery
years, five- and ten-year scarring, and the unemployment increase. Complete-case samples
range from 144 to 178 episodes. Table~\ref{tab:severity} reports the credit-growth
coefficients.

\begin{table}[htbp]
\centering
\small
\caption{Coefficient on three-year credit growth for nine severity outcomes, release 6
annual-decline episodes. Country-clustered standard errors are in parentheses. The Holm
adjustment is across the nine outcomes.}
\label{tab:severity}
\begin{tabular}{lrrrr}
\toprule
Outcome & Coefficient & $p$ & Holm $p$ & Episodes \\
\midrule
Peak-to-trough decline (\%) & 0.093 (0.105) & 0.386 & 1.000 & 178 \\
Five-year loss, trend & $-0.005$ (0.441) & 0.991 & 1.000 & 178 \\
Ten-year loss, trend & $-0.131$ (0.792) & 0.870 & 1.000 & 176 \\
Five-year loss, prior peak & 0.192 (0.358) & 0.599 & 1.000 & 178 \\
Ten-year loss, prior peak & 0.332 (0.494) & 0.511 & 1.000 & 176 \\
Recovery (years) & 0.021 (0.025) & 0.414 & 1.000 & 161 \\
Five-year scarring (\%) & $-0.007$ (0.149) & 0.961 & 1.000 & 178 \\
Ten-year scarring (\%) & $-0.176$ (0.174) & 0.326 & 1.000 & 176 \\
Unemployment increase (pp) & 0.043 (0.046) & 0.370 & 1.000 & 144 \\
\bottomrule
\end{tabular}
\end{table}

No raw $p$-value is below 0.32, and every Holm-adjusted $p$-value is 1.000. Adding
three-year equity returns as a control reduces the samples to 130 to 161 episodes in 16
countries. All nine Holm-adjusted $p$-values for credit growth stay at 1.000. Wide intervals
do not show that the true association is zero. They show that release 6 with this design
cannot distinguish the original effect sizes from zero.

\section{RQ3: Conditional break-even arithmetic}
Let $c$ be an annual output cost of restraint in percent of GDP. Let $\lambda$ be the annual
downturn frequency, and let $\Delta D$ be the discounted output loss avoided per downturn.
Holding $\lambda$ fixed, restraint breaks even in output terms when $c \le \lambda\,\Delta D$.
The required avoided loss is $c/\lambda$. This number is arithmetic by construction. It
depends only on the chosen cost and the counted frequency. It is not an estimate of any
policy effect, and it says nothing about whether any policy achieves it.

The frequency is 0.114 downturns per eligible country-year, counting the 274 downturns with
complete ten-year outcomes. A country-cluster bootstrap with 1,000 draws gives a 95 percent
interval of 0.106 to 0.123. For $c=0.1$ percent, the required avoided loss is 0.87
GDP-percent-years per downturn (interval 0.81 to 0.95). This interval reflects only the
uncertainty in the frequency, because $c$ is fixed by assumption. For $c=0.5$ percent, it is
4.37. The original JST chronology gives a similar number. Counting recessions with complete
five-year windows, it has 207 recessions in 1,442 eligible country-years, and the same
arithmetic gives 0.70 for $c=0.1$ percent. With complete ten-year windows, it has 201
recessions in 1,372 country-years and gives 0.68.

An earlier draft also divided 0.87 by the mean measured ten-year loss per downturn. I do
not report that share as a result, because the measured loss is too large to be credible.
Three features of the loss measure inflate it. First, yearly output gaps are floored at
zero, so years above the counterfactual path never offset years below it. Second, the
ten-year windows overlap: 68 percent of the 299 downturns (66 percent of the 274 with
complete windows) are followed by another downturn in the same country within ten years, so
the same lost year can count more than once. Third, 28 complete episodes start in war years
(1914 to 1918 and 1939 to 1945). At a 3 percent discount rate, their mean loss against the
linear trend is 185 GDP-percent-years, against 76 for the other episodes.

The product $\lambda$ times the mean loss is a sanity check. Under the baseline measure at a
0 percent discount rate, it implies that the average country loses 11.9 percent of GDP every
year against its pre-downturn trend. Table~\ref{tab:sensitivity} recomputes the mean loss
under alternative definitions. Signed gaps remove the floor. De-duplicated years assign each
post-peak year only to the latest downturn whose window covers it, which keeps every episode
and the same frequency. The non-overlapping subset keeps only downturns that start more than
ten years after the previous kept one, so its frequency is lower. The war exclusion drops
war-year starts from both the episodes and the exposure.

\begin{table}[htbp]
\centering
\small
\caption{Sensitivity of the mean ten-year loss per downturn and of the share of it that a
0.1 percent annual cost requires. Trend is the pre-episode linear trend. Peak is the flat
prior-peak path. Annual loss is frequency times mean trend loss, in percent of GDP per year.
A dash marks a negative mean loss, where the share is undefined. Point estimates only.}
\label{tab:sensitivity}
\setlength{\tabcolsep}{4pt}
\begin{tabular}{lrrrrrrr}
\toprule
Loss definition & Rate & $N$ & $c/\lambda$ & Trend loss & Trend share & Peak share & Annual loss \\
\midrule
Baseline (floored, overlapping) & 3\% & 274 & 0.87 & 86.8 & 1.01\% & 3.53\% & 9.9 \\
No floor & 3\% & 274 & 0.87 & 65.7 & 1.33\% & -- & 7.5 \\
De-duplicated, floored & 3\% & 274 & 0.87 & 63.5 & 1.38\% & 3.97\% & 7.3 \\
Non-overlapping, floored & 3\% & 144 & 1.66 & 101.0 & 1.65\% & 6.71\% & 6.1 \\
War starts excluded, floored & 3\% & 246 & 0.89 & 75.6 & 1.17\% & 4.93\% & 8.5 \\
De-duplicated, floored, no war & 3\% & 246 & 0.89 & 50.5 & 1.76\% & 5.92\% & 5.7 \\
No floor, de-duplicated, no war & 3\% & 246 & 0.89 & 39.5 & 2.24\% & -- & 4.5 \\
\midrule
Baseline (floored, overlapping) & 0\% & 274 & 0.87 & 104.0 & 0.84\% & 3.05\% & 11.9 \\
No floor & 0\% & 274 & 0.87 & 78.3 & 1.12\% & -- & 9.0 \\
De-duplicated, floored & 0\% & 274 & 0.87 & 74.6 & 1.17\% & 3.47\% & 8.5 \\
Non-overlapping, floored & 0\% & 144 & 1.66 & 121.0 & 1.37\% & 5.84\% & 7.3 \\
War starts excluded, floored & 0\% & 246 & 0.89 & 90.8 & 0.98\% & 4.26\% & 10.2 \\
De-duplicated, floored, no war & 0\% & 246 & 0.89 & 59.1 & 1.50\% & 5.20\% & 6.7 \\
No floor, de-duplicated, no war & 0\% & 246 & 0.89 & 46.3 & 1.92\% & -- & 5.2 \\
\bottomrule
\end{tabular}
\end{table}

Each correction lowers the trend loss, but the implied annual trend loss stays between 4.5
and 11.9 percent of GDP. That is still implausibly large, which points to the five-year
linear trend as a further source of bias. Against the flat prior peak, signed gaps give a
negative mean loss: on average, output over the ten years after the peak is above the prior
peak.
The required share therefore ranges from 0.84 to 6.71 percent across floored definitions
and is undefined for signed gaps against the prior peak. Because the share depends on the
loss definition, I report only the 0.87 arithmetic as a result.

\section{Limits}
The reproduction checks the archived \texttt{table8.tex} in the authors' replication archive,
in Python, not in Stata. I did not compare it with the printed 2013 journal table. The
fragility results use 14 to 18 country clusters, which makes clustered inference fragile.
The vintage comparison cannot separate data revision from a possible change in the credit
definition, and it has no formal test of the difference. The dating comparison cannot
separate the dating rule from sample composition. All estimates are conditional
associations. Selecting on output contractions can bias them. The severity outcomes in
Table~\ref{tab:severity} use the same floored and overlapping loss windows as the
break-even section.

The break-even number is arithmetic. It does not estimate $D(m)$, the effect of a policy on
downturn loss, as \citet{svensson2017} and \citet{bgnn2020} do for crisis outcomes and
growth-at-risk. A user of the threshold must supply that estimate from other evidence, and
must also supply a credible counterfactual loss per downturn.

\section{Conclusion}
An independent implementation of the aggregate \citet{jst2013} Table 8 matches the authors'
archived output at three decimals. On release 6 data, the credit-severity association is
imprecise. Its size changes with the data vintage. Its sign changes only with the
subsample, the control set, or the episode rule. No severity coefficient survives a
Holm adjustment. By construction, a 0.1 percent annual cost requires 0.87 GDP-percent-years
of avoided loss per downturn. The measured loss per downturn is too inflated to turn that
number into a credible share. Whether any policy achieves the reduction is a separate
question that this note does not answer.

\section*{Data and code availability}
The note does not redistribute third-party data. \texttt{experiments/acquire\_data.py}
downloads the JST release 6 file (CC BY-NC-SA 4.0) and the authors' replication archive
(no license stated). It checks each file against a recorded SHA-256 hash and writes
\texttt{data/source\_manifest.json}. The analysis scripts, run records, and the README
commands are in the project folder. Each run folder records its command, inputs, code
hashes, outputs, and environment.

\begin{thebibliography}{30}
\bibitem[Baron et~al.(2021)]{baron2021}
Baron, M., Verner, E., and Xiong, W. (2021).
\newblock Banking crises without panics.
\newblock \emph{Quarterly Journal of Economics}, 136(1), 51--113.
\newblock \url{https://doi.org/10.1093/qje/qjaa034}.

\bibitem[Bordo and Haubrich(2017)]{bordohaubrich2017}
Bordo, M.~D., and Haubrich, J.~G. (2017).
\newblock Deep recessions, fast recoveries, and financial crises: Evidence from the American
  record.
\newblock \emph{Economic Inquiry}, 55(1), 527--541.
\newblock \url{https://doi.org/10.1111/ecin.12374}.

\bibitem[Bridges et~al.(2017)]{bridges2017}
Bridges, J., Jackson, C., and McGregor, D. (2017).
\newblock Down in the slumps: The role of credit in five decades of recessions.
\newblock Bank of England Staff Working Paper No. 659.
\newblock \url{https://www.bankofengland.co.uk/working-paper/2017/down-in-the-slumps-the-role-of-credit-in-five-decades-of-recessions}.

\bibitem[Brand\~ao-Marques et~al.(2020)]{bgnn2020}
Brand\~ao-Marques, L., Gelos, R.~G., Narita, M., and Nier, E. (2020).
\newblock Leaning against the wind: A cost-benefit analysis for an integrated policy
  framework.
\newblock IMF Working Paper WP/20/123, International Monetary Fund.
\newblock \url{https://doi.org/10.5089/9781513549651.001}.

\bibitem[Claessens et~al.(2012)]{claessens2012}
Claessens, S., Kose, M.~A., and Terrones, M.~E. (2012).
\newblock How do business and financial cycles interact?
\newblock \emph{Journal of International Economics}, 87(1), 178--190.
\newblock \url{https://doi.org/10.1016/j.jinteco.2011.11.008}.

\bibitem[Gourinchas and Obstfeld(2012)]{gourinchasobstfeld2012}
Gourinchas, P.-O., and Obstfeld, M. (2012).
\newblock Stories of the twentieth century for the twenty-first.
\newblock \emph{American Economic Journal: Macroeconomics}, 4(1), 226--265.
\newblock \url{https://doi.org/10.1257/mac.4.1.226}.

\bibitem[Greenwood et~al.(2022)]{greenwood2022}
Greenwood, R., Hanson, S.~G., Shleifer, A., and S{\o}rensen, J.~A. (2022).
\newblock Predictable financial crises.
\newblock \emph{Journal of Finance}, 77(2), 863--921.
\newblock \url{https://doi.org/10.1111/jofi.13105}.

\bibitem[Jord\`a et~al.(2022)]{jkst2022}
Jord\`a, \`O., Kornejew, M., Schularick, M., and Taylor, A.~M. (2022).
\newblock Zombies at large? Corporate debt overhang and the macroeconomy.
\newblock \emph{Review of Financial Studies}, 35(10), 4561--4586.
\newblock \url{https://doi.org/10.1093/rfs/hhac018}.

\bibitem[Jord\`a et~al.(2015a)]{jst2015bubbles}
Jord\`a, \`O., Schularick, M., and Taylor, A.~M. (2015).
\newblock Leveraged bubbles.
\newblock \emph{Journal of Monetary Economics}, 76(Supplement), S1--S20.
\newblock \url{https://doi.org/10.1016/j.jmoneco.2015.08.005}.

\bibitem[Jord\`a et~al.(2016a)]{jst2016mortgaging}
Jord\`a, \`O., Schularick, M., and Taylor, A.~M. (2016).
\newblock The great mortgaging: Housing finance, crises and business cycles.
\newblock \emph{Economic Policy}, 31(85), 107--152.
\newblock \url{https://doi.org/10.1093/epolic/eiv017}.

\bibitem[Jord\`a et~al.(2016b)]{jst2016sovereigns}
Jord\`a, \`O., Schularick, M., and Taylor, A.~M. (2016).
\newblock Sovereigns versus banks: Credit, crises, and consequences.
\newblock \emph{Journal of the European Economic Association}, 14(1), 45--79.
\newblock \url{https://doi.org/10.1111/jeea.12144}.

\bibitem[Jord\`a et~al.(2013)]{jst2013}
Jord\`a, \`O., Schularick, M., and Taylor, A.~M. (2013).
\newblock When credit bites back.
\newblock \emph{Journal of Money, Credit and Banking}, 45(S2), 3--28.
\newblock \url{https://doi.org/10.1111/jmcb.12069}.

\bibitem[Jord\`a et~al.(2017)]{jstdatabase}
Jord\`a, \`O., Schularick, M., and Taylor, A.~M. (2017).
\newblock Macrofinancial history and the new business cycle facts.
\newblock \emph{NBER Macroeconomics Annual}, 31, 213--263.
\newblock \url{https://doi.org/10.1086/690241}. Database release 6, CC BY-NC-SA 4.0.

\bibitem[Jord\`a et~al.(2015b)]{jstcode}
Jord\`a, \`O., Schularick, M., and Taylor, A.~M. (2015).
\newblock ``When Credit Bites Back'' replication archive, version of 14 August 2015.
\newblock \url{https://www.moritzschularick.com/academic/datasets-and-codes/}.

\bibitem[Krishnamurthy and Muir(2025)]{krishnamurthymuir2025}
Krishnamurthy, A., and Muir, T. (2025).
\newblock How credit cycles across a financial crisis.
\newblock \emph{Journal of Finance}, 80(3), 1339--1378.
\newblock \url{https://doi.org/10.1111/jofi.13431}.

\bibitem[Mian et~al.(2017)]{msv2017}
Mian, A., Sufi, A., and Verner, E. (2017).
\newblock Household debt and business cycles worldwide.
\newblock \emph{Quarterly Journal of Economics}, 132(4), 1755--1817.
\newblock \url{https://doi.org/10.1093/qje/qjx017}.

\bibitem[M\"uller and Verner(2024)]{mullerverner2024}
M\"uller, K., and Verner, E. (2024).
\newblock Credit allocation and macroeconomic fluctuations.
\newblock \emph{Review of Economic Studies}, 91(6), 3645--3676.
\newblock \url{https://doi.org/10.1093/restud/rdad112}.

\bibitem[Richter et~al.(2021)]{rsw2021}
Richter, B., Schularick, M., and Wachtel, P. (2021).
\newblock When to lean against the wind.
\newblock \emph{Journal of Money, Credit and Banking}, 53(1), 5--39.
\newblock \url{https://doi.org/10.1111/jmcb.12701}.

\bibitem[Romer and Romer(2017)]{romerromer2017}
Romer, C.~D., and Romer, D.~H. (2017).
\newblock New evidence on the aftermath of financial crises in advanced countries.
\newblock \emph{American Economic Review}, 107(10), 3072--3118.
\newblock \url{https://doi.org/10.1257/aer.20150320}.

\bibitem[Schularick and Taylor(2012)]{schularicktaylor2012}
Schularick, M., and Taylor, A.~M. (2012).
\newblock Credit booms gone bust: Monetary policy, leverage cycles, and financial crises,
  1870--2008.
\newblock \emph{American Economic Review}, 102(2), 1029--1061.
\newblock \url{https://doi.org/10.1257/aer.102.2.1029}.

\bibitem[Sufi and Taylor(2022)]{sufitaylor2022}
Sufi, A., and Taylor, A.~M. (2022).
\newblock Financial crises: A survey.
\newblock In \emph{Handbook of International Economics}, Vol.~6, 291--340. Elsevier.
\newblock \url{https://doi.org/10.1016/bs.hesint.2022.02.012}.

\bibitem[Summers(2017)]{summers2017}
Summers, P.~M. (2017).
\newblock Credit booms gone bust: Replication of {S}chularick and {T}aylor ({AER} 2012).
\newblock \emph{Journal of Applied Econometrics}, 32(5), 1033--1038.
\newblock \url{https://doi.org/10.1002/jae.2554}.

\bibitem[Svensson(2017)]{svensson2017}
Svensson, L.~E.~O. (2017).
\newblock Cost-benefit analysis of leaning against the wind.
\newblock \emph{Journal of Monetary Economics}, 90, 193--213.
\newblock \url{https://doi.org/10.1016/j.jmoneco.2017.07.004}.
\end{thebibliography}
\end{document}
